
\begingroup
\shorthandoff{<>}% Babel español trata de otro modo ciertos delimitadores de Xy-pic.
\subsection{}
\label{XIII.2.2}
Conservemos las notaciones de~\Ref{XIII.2.1}. Vamos a definir una noción de
objeto moderadamente ramificado de un stack~$\Phi$ sobre~$U$ cuando este viene
dado, localmente \emph{sobre}~$X$ y sobre~$S$ para la topología étale, como
\emph{imagen inversa de un stack}~$\Psi$ \emph{sobre}~$S$.

Sea primero $G$ un gerbe sobre~$U$ y supongamos dados un morfismo étale
sobreyectivo $S_1\to S$, un morfismo étale sobreyectivo
$X_2\to X\times_S S_1$, un gerbe trivial $H$ sobre~$S_1$ y un isomorfismo
$$
G|{U_2}\to H|{U_2},
$$
donde
\marginpar{377}
\ifthenelse{\boolean{orig}}
{$U_1=U\times_X X_1$, $U_2=U\times_X X_2$.}
{$X_1=X\times_S S_1$, $U_1=U\times_X X_1$,
$U_2=U\times_X X_2$.}
Al elegir una trivialización de~$H|{X_2}$, el isomorfismo anterior identifica
$G|{U_2}$ con el stack de torsores bajo un haz de grupos~$F$. Se dice que un
objeto $x$ de~$G_U$ está moderadamente ramificado sobre~$X$ relativamente
\ifthenelse{\boolean{orig}}
{a~$S$,}
{a~$S$}
si la restricción de~$x$ a~$U_2$ es un torsor moderadamente ramificado
\ifthenelse{\boolean{orig}}
{sobre~$X$ relativamente a~$S$.}
{sobre~$X_2$ relativamente a~$S_1$.}
Según~\Ref{XIII.2.1.4}, esta noción no depende de la trivialización elegida
de~$H|X_2$.

Sea ahora $\Phi$ un stack sobre~$U$ y supongamos dados un morfismo étale
sobreyectivo $S_1\to S$, un morfismo étale sobreyectivo
$X_2\to X\times_S S_1$, un stack $\Psi$ sobre~$S_1$ y un isomorfismo
$$
i\colon \Phi|U_2\to \Psi|U_2.
$$
Sea $x$ un objeto de~$\Phi_U$, sea $G_x$ el subgerbe maximal de~$\Phi$
generado por~$x$
\ifthenelse{\boolean{orig}}
{[III~2.1.7]}
{\cite[III~2.1.7]{XIII.1}},
y $S\Phi$ el haz de los subgerbes maximales de~$\Phi$. El isomorfismo~$i$
induce un isomorfismo
$$
S\Phi|U_2\to S\Psi|U_2.
$$
De~\Ref{XIII.5.7} resulta que, tras reemplazar si es preciso $S_1$ por una
extensión étale sobreyectiva, existe un único subgerbe maximal~$H$ de~$\Psi$,
que puede suponerse trivial, tal que~$i$ define un isomorfismo
$$
G_x|U_2\to H|U_2.
$$
Se dice que el objeto~$x$ está moderadamente ramificado sobre~$X$
relativamente a~$S$ si lo está como objeto de~$G_x$, provisto del
isomorfismo anterior.

\subsubsection{}
\label{XIII.2.2.1}
Sea $\Phi$ un stack sobre~$U$ dado, localmente sobre~$X$ y~$S$, como imagen
inversa de un stack sobre~$S$, y sea
$$
\xymatrix{ U \ar[r]^i
\ar[d]_g & X \ar[dl]^f \\ T }
$$
un diagrama como en~\Ref{XIII.2.1.2}. Para todo esquema $T'$ étale sobre~$T$,
\ifthenelse{\boolean{orig}}
{si $U'=U\times_TT'$,}
{si $U'=U\times_TT'$ y $X'=X\times_TT'$,}
consideremos el subconjunto
$(g_*^{\tame}\Phi)_{T'}$
\label{indnot:mg}\oldindexnot{$(g_*^{\tame}\Phi)_{T'}$|hyperpage}%
de~$(g_*\Phi)_{T'}=\Phi_{U'}$ formado
\marginpar{378}
por los objetos de~$\Phi_{U'}$ moderadamente ramificados
\index{objeto moderadamente ramificado de un stack|hyperpage}%
\ifthenelse{\boolean{orig}}
{sobre~$X$ relativamente a~$S$.}
{sobre~$X'$ relativamente a~$S$.}
Se denomina \emph{imagen directa
moderadamente ramificada}
\index{imagen directa moderadamente ramificada de un stack|hyperpage}%
\index{moderadamente ramificada (imagen directa de un stack)|hyperpage}%
de~$\Phi$ por~$g$, y se denota por
$$
g_*^{\tame}\Phi
$$
la subcategoría plena de~$g_*\Phi$ cuyos objetos sobre un esquema $T'$ étale
sobre~$T$ son los objetos de~$(g_*^{\tame}\Phi)_{T'}$. Es claro que
$g_*^{\tame}\Phi$ es un substack de~$g_*\Phi$.

\subsubsection{}
\label{XIII.2.2.2}
\ifthenelse{\boolean{orig}}
{}
{Retomemos las notaciones de~\Ref{XIII.2.1.5}. }
Reduciendo al caso de un stack de torsores, se ve que, si $\Phi$ es un stack
sobre~$U$ que, localmente para la topología étale sobre~$S$ y sobre~$X$, es
imagen inversa de un stack sobre~$S$, el morfismo canónico
$$
h^*(g_*\Phi)\to g'_*\Phi'
$$
induce por restricción un morfismo canónico
$$
h^*(g_*^{\tame}\Phi)\to g_*^{\prime\tame}\Phi'
$$

\setcounter{proposition}{2}% La observación siguiente es XIII.2.3; los §§2.0--2.2
\begin{remarques}
\label{XIII.2.3}

a) Si $F$ es un haz de conjuntos localmente constante y constructible
sobre~$U$, para que $F$ esté moderadamente ramificado sobre~$X$ relativamente
a~$S$, basta que la condición de~\Ref{XIII.2.1.1} se satisfaga para los puntos
geométricos de~$S$ situados sobre los puntos maximales de~$S$. Para verlo,
puede suponerse que $D$ es de cruzamientos normales estrictos. El haz~$F$ es
representable por un recubrimiento étale~$V$ de~$U$. Si $\overline{s}$ es un
punto geométrico de~$S$ y $y$ un punto maximal de~$Y_{\overline{s}}$,
\ifthenelse{\boolean{orig}}
{denotemos por $\overline{S}$ la localización estricta de~$S$ en
$\overline{s}$ y por $\overline{X}$ la localización estricta de~$X$ en
$\overline{y}$,}
{elijamos un punto geométrico $\overline y$ situado sobre~$y$; denotemos por
$\overline{S}$ la localización estricta de~$S$ en~$\overline{s}$ y por
$\overline{X}$ la localización estricta de~$X$ en~$\overline{y}$,}
$\overline{U}=U\times_X\overline{X}$ y
$\overline{V}=V\times_X\overline{X}$. Si la condición
de~\Ref{XIII.2.1.1} se satisface en los puntos geométricos situados sobre los
puntos maximales de~$S$, de~\Ref{XIII.5.5} más abajo resulta que
$\overline{V}$ es un recubrimiento de~$\overline{U}$ moderadamente ramificado
sobre~$\overline{X}$ relativamente a~$\overline{S}$; por tanto, $V$ es un
recubrimiento étale de~$U$ moderadamente ramificado sobre~$X$ relativamente
a~$S$.

b)
\marginpar{379}
Sea $F$ un haz de grupos sobre~$U$ moderadamente ramificado sobre~$X$
relativamente a~$S$. Si $\overline{s}$ es un punto geométrico de~$S$ y $y$
un punto maximal de~$Y_{\overline{s}}$, denotemos por $K$ el cuerpo de
fracciones de~$\cal{O}_{X_{\overline{s}},y}$. Supongamos que, para todo punto
$\overline{s}$ y todo punto~$y$, la $K$-álgebra~$L$ cuyo espectro representa
$F|K$ es de rango primo con la característica residual~$p$ de
\ifthenelse{\boolean{orig}}
{$\cal{O}_{X_{\overline{s}}}$.}
{$\cal{O}_{X_{\overline{s}},y}$.}
Diremos a veces, por abuso de lenguaje, que $F$ es primo con las
características residuales de~$S$.
\index{primo con las características residuales (haz)|hyperpage}%
En tal caso, todo torsor~$P$ bajo~$F$ está moderadamente ramificado sobre~$X$
relativamente a~$S$.

En efecto,
\ifthenelse{\boolean{orig}}
{sean}
{elijamos un punto geométrico $\overline y$ situado sobre~$y$; sean}
$\overline{R}$ la localización estricta de
$\cal{O}_{X_{\overline{s}},y}$ en~$\overline{y}$, $\overline{K}$ su cuerpo de
fracciones y $\overline{F}$ la imagen inversa de~$F$ sobre~$\overline{K}$.
Demostremos que puede suponerse constante
\ifthenelse{\boolean{orig}}
{$F$.}
{$\overline{F}$.}
Como $F$ está moderadamente ramificado sobre~$X$ relativamente a~$S$,
$\overline{F}$ es representable por el espectro de una $\overline{K}$-álgebra
$L=\prod L_i$, donde los $L_i$ son extensiones de~$\overline{K}$ de grado
primo con~$p$. Puede encontrarse, por tanto, una extensión~$K'$ de
$\overline{K}$ de grado primo con~$p$ tal que $\overline{F}|K'$ sea un haz
constante. Según~\Ref{XIII.2.0.3}, para demostrar que $P|\overline{K}$ está
moderadamente ramificado sobre~$\overline{R}$, basta ver que $P|K'$ está
moderadamente ramificado sobre la clausura integral de~$\overline{R}$ en
$K'$; esto reduce al caso en que $\overline{F}$ es constante. Supongamos en
adelante que $\overline{F}$ es constante. La $\overline{K}$-álgebra~$H$ que
representa $P|\overline{K}$ es entonces producto de extensiones~$H_i$ de
$\overline{K}$ isomorfas entre sí. Como el rango de~$H$ es primo con~$p$,
también lo es $[H_1:\overline{K}]$, lo que demuestra que $H$ está
moderadamente ramificada
\ifthenelse{\boolean{orig}}
{sobre~$X$ relativamente a~$S$.}
{sobre~$\overline{R}$.}

c) Sean $X$ un esquema regular, $D$ un divisor de cruzamientos normales
sobre~$X$ (SGA~5 I~3.1.5), $U=X-\Supp D$ y $F$ un haz de conjuntos sobre~$U$.
Si $y$ es un punto maximal de~$\Supp D$, denotemos por $K$ el cuerpo de
fracciones de~$\cal{O}_{X,y}$. Se dice que $F$ está
\emph{moderadamente ramificado relativamente a}~$D$
\index{moderadamente ramificado (haz)|hyperpage}%
si, para todo punto maximal $y$ de~$\Supp D$, $F|K$ es representable por el
espectro de una $K$-álgebra moderadamente ramificada sobre~$\cal{O}_{X,y}$.
\end{remarques}
\endgroup
